
\section{滤过 $A_\infty$ 代数的代数框架}
\label{sec:3-2}

本节将在纯代数语境中引入滤过 $A_\infty$ 代数和 $A_\infty$ 同态的概念，
因而不涉及拉格朗日子流形。我们的 $A_\infty$ 代数定义与文献中的标准定义
（\textbf{[Sta63]}）略有不同，因为这里的系数环是滤过交换环
$\Lambda_{0,\mathrm{nov}}$。为此，第 3.2.1 小节先回顾无滤过
$A_\infty$ 代数及其同态，下一小节再引入滤过 $A_\infty$ 代数和滤过
$A_\infty$ 同态。为了把通常的无滤过 $A_\infty$ 代数同这里的滤过
$A_\infty$ 代数区分开来，我们在前者的记号上加横线。

\subsection{$A_\infty$ 代数与同态}
\label{subsec:3-2-1}

令 $R$ 为带幺交换环，并令
$\overline C=\bigoplus_{m\in\mathbb Z}\overline C^m$ 为自由分次
$R$-模。我们始终假定 $\overline C$ 由可数多个元素生成。置
$(\overline C[1])^m=\overline C^{m+1}$，并令
\[
  B_k(\overline C[1])
  =\bigoplus_{m_1,\ldots,m_k}
  (\overline C[1])^{m_1}\otimes\cdots\otimes
  (\overline C[1])^{m_k}.
\]
有些文献把上式右端记作 $B_k(\overline C)$，而不是
$B_k(\overline C[1])$。不过，本书始终采用后一个记号；有时也把
$B_k(\overline C[1])$ 简写为 $B_k\overline C[1]$。假设给定一列次数为
$+1$ 的映射 $\overline{\mathfrak m}
=\{\overline{\mathfrak m}_k\}_{k\geq1}$：
\[
  \overline{\mathfrak m}_k\colon B_k(\overline C[1])
  \longrightarrow\overline C[1],
  \qquad k=1,2,\ldots.
\]

直和
$B(\overline C[1]):=\bigoplus_k B_k(\overline C[1])$ 具有分次余代数结构。
其余积\liftIndex{I035} $\Delta$ 定义为
\[
  \Delta(x_1\otimes\cdots\otimes x_n)
  =\sum_{k=0}^{n}
  (x_1\otimes\cdots\otimes x_k)
  \otimes(x_{k+1}\otimes\cdots\otimes x_n).
\]
可以把 $\overline{\mathfrak m}_k$ 唯一延拓为余导子\liftIndex{I092}
\[
  \widehat{\overline{\mathfrak m}}_k\colon
  \bigoplus_n B_n(\overline C[1])
  \longrightarrow
  \bigoplus_n B_{n-k+1}(\overline C[1]),
\]
其公式为
\begin{equation}
\begin{aligned}
  &\widehat{\overline{\mathfrak m}}_k
  (x_1\otimes\cdots\otimes x_n)\\
  &\quad={}
  \sum_{\ell=1}^{n-k+1}
  (-1)^{\deg x_1+\cdots+\deg x_{\ell-1}+\ell-1}
  x_1\otimes\cdots\otimes
  \overline{\mathfrak m}_k
  (x_\ell,\ldots,x_{\ell+k-1})
  \otimes\cdots\otimes x_n .
\end{aligned}
\tag{3.2.1}\label{eq:3-2-1}
\end{equation}
这里 $k\leq n$；当 $k>n$ 时，令
$\widehat{\overline{\mathfrak m}}_k=0$。这里以及下文中，$\deg x$ 都指
移位之前 $x$ 的次数。

以后，对 $x\in(\overline C[1])^m=\overline C^{m+1}$，置
\begin{equation}
  \deg' x=m=\deg x-1.
  \tag{3.2.2}\label{eq:3-2-2}
\end{equation}
这称为移位次数\liftIndex{I028}。用这个记号，式~\eqref{eq:3-2-1} 中的符号就是
$(-1)^{\deg' x_1+\cdots+\deg' x_{\ell-1}}$。置
\[
  \widehat{\overline d}
  =\sum_k\widehat{\overline{\mathfrak m}}_k.
\]

\begin{definition}{3.2.3}\label{def:3-2-3}
若 $\widehat{\overline d}\circ\widehat{\overline d}=0$，就称
$\overline{\mathfrak m}=\{\overline{\mathfrak m}_k\}_{k\geq1}$
在 $\overline C$ 上定义了一个 $R$ 上的
$A_\infty$ 代数结构\liftIndex{I002}。
\end{definition}

\begin{remark}{3.2.4}\label{rem:3-2-4}
\begin{enumerate}[label=\textup{(\arabic*)},leftmargin=2.8em]
  \item 算子 $\overline{\mathfrak m}_k$ 只对 $k>0$ 定义，也就是说
  $\overline{\mathfrak m}_0=0$。对下一小节定义的滤过 $A_\infty$ 代数，
  则要考虑所有 $k\geq0$ 的 $\mathfrak m_k$。由于
  $\overline{\mathfrak m}_0=0$，有
  $\overline{\mathfrak m}_1\overline{\mathfrak m}_1=0$，所以得到上链复形
  $(\overline C[1],\overline{\mathfrak m}_1)$；另见注 3.2.21 (1)。

  \item 把 $R$ 上的微分分次代数（D.G.A.）视为 $A_\infty$ 代数时，只需令
  $\overline{\mathfrak m}_k=0$（$k\geq3$）。不过要注意，
  $\overline{\mathfrak m}_2$ 的符号与 D.G.A. 中乘积的符号不同：
  \begin{equation}
  \begin{aligned}
    \overline{\mathfrak m}_1(x)
      &=(-1)^{\deg x}dx,\\
    \overline{\mathfrak m}_2(x,y)
      &=(-1)^{\deg x(\deg y+1)}x\mathbin{\cdot}y.
  \end{aligned}
  \tag{3.2.5}\label{eq:3-2-5}
  \end{equation}
  这里 $x\mathbin{\cdot}y$ 是 D.G.A. 中的乘积；另见下面的注 3.2.21 (2)。
\end{enumerate}
\end{remark}

下面定义 $A_\infty$ 同态。令
$(\overline C_i,\overline{\mathfrak m}^{\,i})$（$i=1,2$）为带幺交换环
$R$ 上的 $A_\infty$ 代数。对 $k=1,2,\ldots$，考虑次数为 $0$ 的映射族
\[
  \overline f_k\colon B_k(\overline C_1[1])
  \longrightarrow\overline C_2[1].
\]
这些映射诱导
$\widehat{\overline f}\colon B(\overline C_1[1])\to
B(\overline C_2[1])$；它在 $B_k(\overline C_1[1])$ 上由下式给出：
\begin{equation}
\begin{aligned}
  &\widehat{\overline f}(x_1\otimes\cdots\otimes x_k)\\
  &\quad={}
  \sum_{0<k_1<\cdots<k_n<k}
  \overline f_{k_1}(x_1,\ldots,x_{k_1})\otimes\cdots\\
  &\qquad\quad{}\otimes
  \overline f_{k_{i+1}-k_i}
  (x_{k_i+1},\ldots,x_{k_{i+1}})\otimes\cdots\\
  &\qquad\quad{}\otimes
  \overline f_{k-k_n}(x_{k_n+1},\ldots,x_k).
\end{aligned}
\tag{3.2.6}\label{eq:3-2-6}
\end{equation}
容易看出，$\widehat{\overline f}$ 是一个余代数同态。这里
$\overline f_0=0$。

\begin{definition}{3.2.7}\label{def:3-2-7}
若
\[
  \widehat{\overline f}\circ\widehat{\overline d}^{\,1}
  =\widehat{\overline d}^{\,2}\circ\widehat{\overline f},
\]
就称 $\overline f=\{\overline f_k\}_{k\geq1}$ 是从
$\overline C_1$ 到 $\overline C_2$ 的
$A_\infty$ 同态\liftIndex{I087}，简记为
$\overline f\colon\overline C_1\to\overline C_2$。
\end{definition}

根据定义，从 $\overline C_1$ 到 $\overline C_2$ 的 $A_\infty$ 同态
$\overline f$ 诱导同态
\[
  H^*\bigl(B(\overline C_1[1]),\widehat{\overline d}^{\,1}\bigr)
  \longrightarrow
  H^*\bigl(B(\overline C_2[1]),\widehat{\overline d}^{\,2}\bigr).
\]
设
$\overline f_k^{\,i}\colon B_k(\overline C_i[1])\to
\overline C_{i+1}[1]$（$i=1,2$）各自定义一个 $A_\infty$ 同态。

\begin{definition}{3.2.8}\label{def:3-2-8}
$\overline f^{\,1}$ 与 $\overline f^{\,2}$ 的复合
$\overline f^{\,2}\circ\overline f^{\,1}
=\{(\overline f^{\,2}\circ\overline f^{\,1})_k\}_{k\geq1}$
\liftIndex{I088}由下式定义：
\[
\begin{aligned}
  &(\overline f^{\,2}\circ\overline f^{\,1})_k
  (x_1,\ldots,x_k)\\
  &\quad={}
  \sum_m\sum_{k_1+\cdots+k_m=k}
  \overline f_m^{\,2}\bigl(
    \overline f_{k_1}^{\,1}(x_1,\ldots,x_{k_1}),\ldots,
    \overline f_{k_m}^{\,1}(x_{k-k_m+1},\ldots,x_k)
  \bigr).
\end{aligned}
\]
\end{definition}

\begin{lemma}{3.2.9}\label{lem:3-2-9}
$\overline f^{\,2}\circ\overline f^{\,1}$ 定义了从
$\overline C_1$ 到 $\overline C_3$ 的 $A_\infty$ 同态。
\end{lemma}

\begin{proof}
若 $\widehat{\overline f}^{\,1}$、$\widehat{\overline f}^{\,2}$ 分别由
$\overline f_k^{\,1}$、$\overline f_k^{\,2}$ 通过~\eqref{eq:3-2-6} 诱导，
则复合 $\widehat{\overline f}^{\,2}\circ
\widehat{\overline f}^{\,1}$ 由
$(\overline f^{\,2}\circ\overline f^{\,1})_k$ 诱导，结论立即成立。
\end{proof}

注意
\[
  \widehat{\overline f}^{\,2}\circ\widehat{\overline f}^{\,1}
  =\widehat{\overline f^{\,2}\circ\overline f^{\,1}}.
\]
次数为 $+1$ 的映射 $\overline{\mathfrak m}_1$ 满足
$\overline{\mathfrak m}_1\circ\overline{\mathfrak m}_1=0$（见注 3.2.4
(1)），因而有复形 $(\overline C[1],\overline{\mathfrak m}_1)$。
$A_\infty$ 同态会诱导上同调之间的同态。

\begin{definition}{3.2.10}\label{def:3-2-10}
令 $(\overline C_i,\overline{\mathfrak m}^{\,i})$ 为 $R$ 上的
$A_\infty$ 代数。若 $A_\infty$ 同态
\[
  \overline f\colon
  (\overline C_1,\overline{\mathfrak m}^{\,1})
  \longrightarrow
  (\overline C_2,\overline{\mathfrak m}^{\,2})
\]
的线性项 $\overline f_1$ 诱导上链同伦等价，就称 $\overline f$ 为
\emph{弱同伦等价}。
\end{definition}

\begin{remark}{3.2.11}\label{rem:3-2-11}
\begin{enumerate}[label=\textup{(\arabic*)},leftmargin=2.8em]
  \item 当 $R$ 是域或 $\mathbb Z$ 时，$\overline f$ 是弱同伦等价，当且仅当
  $\overline f_1$ 是拟同构；也就是说，$\overline f_1$ 诱导同构
  \[
    H^*(\overline C_1[1],\overline{\mathfrak m}^{\,1}_1)
    \xrightarrow{\ \cong\ }
    H^*(\overline C_2[1],\overline{\mathfrak m}^{\,2}_1).
  \]
  回忆我们假定 $\overline C_i$ 都是由可数多个元素生成的自由 $R$-模。

  \item 第 4 章将说明，弱同伦等价在 $A_\infty$ 代数之间定义了一个等价关系；
  见推论 4.2.44 和定理 4.2.45。
\end{enumerate}
\end{remark}

\subsection{滤过 $A_\infty$ 代数与同态}
\label{subsec:3-2-2}

现在引入滤过 $A_\infty$ 代数的概念。令
$\bigoplus_{m\in\mathbb Z}C^m$ 为自由分次
$\Lambda_{0,\mathrm{nov}}$-模。对每个 $\lambda\in\mathbb R_{\geq0}$，
$C^m$ 上有滤过 $F^\lambda C^m$，满足
\begin{align}
  F^\lambda C^m&\subset F^{\lambda'}C^m
    \quad\text{若 }\lambda>\lambda',
    \tag{3.2.12.1}\label{eq:3-2-12-1}\\
  T^{\lambda_0}\cdot F^\lambda C^m
    &\subset F^{\lambda+\lambda_0}C^m,
    \tag{3.2.12.2}\label{eq:3-2-12-2}\\
  e^kC^m&\subset C^{m+2k},
    \tag{3.2.12.3}\label{eq:3-2-12-3}\\
  C^m&\text{ 关于该滤过完备},
    \tag{3.2.12.4}\label{eq:3-2-12-4}\\
  C^m&\text{ 有一组基 }\{v_i\}\text{，使得 }
  v_i\in F^0C^m\setminus\bigcup_{\lambda>0}F^\lambda C^m.
    \tag{3.2.12.5}\label{eq:3-2-12-5}
\end{align}
称滤过 $F C$ 为能量滤过，并以 $C$ 表示
$\bigoplus_{m\in\mathbb Z}C^m$ 关于能量滤过的完备化。式
\eqref{eq:3-2-12-3} 表示 $e$ 的次数为 $2$。考虑由
$(C[1])^m=C^{m+1}$ 定义的移位复形 $C[1]$，并令
\[
  B_k(C[1])
  =\bigoplus_{m_1,\ldots,m_k}
  (C[1])^{m_1}\otimes\cdots\otimes(C[1])^{m_k}.
\]
假设给定一列次数为 $+1$ 的映射
$\mathfrak m=\{\mathfrak m_k\}_{k\geq0}$：
\[
  \mathfrak m_k\colon B_k(C[1])\longrightarrow C[1],
  \qquad k=0,1,\ldots.
\]
注意 $\mathfrak m_0\colon\Lambda_{0,\mathrm{nov}}\to C[1]$。我们假定
\begin{align}
  &\mathfrak m_k\bigl(
  F^{\lambda_1}C^{m_1}\otimes\cdots\otimes
  F^{\lambda_k}C^{m_k}\bigr)
  \subset
  F^{\lambda_1+\cdots+\lambda_k}
  C^{m_1+\cdots+m_k-k+2},
  \tag{3.2.12.6}\label{eq:3-2-12-6}\\
  &\mathfrak m_0(1)\in F^{\lambda'}C[1]
  \qquad\text{对某个 }\lambda'>0.
  \tag{3.2.12.7}\label{eq:3-2-12-7}
\end{align}
可见，对滤过 $A_\infty$ 代数，运算 $\mathfrak m_k$ 也在 $k=0$ 时定义。

由于 $C$ 是自由 $\Lambda_{0,\mathrm{nov}}$-模，存在自由 $R$-模
$\overline C$ 以及同构
\begin{equation}
  \overline C\otimes_R\Lambda_{0,\mathrm{nov}}\cong C.
  \tag{3.2.13.1}\label{eq:3-2-13-1}
\end{equation}
此外还假定 $\overline C$ 具有下面的性质。环 $\Lambda_{0,\mathrm{nov}}$
有理想 $\Lambda^+_{0,\mathrm{nov}}$，并且
\[
  \Lambda_{0,\mathrm{nov}}/\Lambda^+_{0,\mathrm{nov}}
  \cong R[e,e^{-1}]
\]
（见 (Conv.4) 与 (1.37)）。对 $k=1,2,\ldots$，置
\begin{equation}
  \overline{\mathfrak m}_k(x_1,\ldots,x_k)
  \equiv\mathfrak m_k(x_1,\ldots,x_k)
  \pmod{\Lambda^+_{0,\mathrm{nov}}C}.
  \tag{3.2.13.2}\label{eq:3-2-13-2}
\end{equation}
这里 $\overline{\mathfrak m}_k$ 是同态
\[
  B_k\overline C[1]\otimes_RR[e,e^{-1}]
  \longrightarrow\overline C[1]\otimes_RR[e,e^{-1}].
\]
我们假定它由一个 $R$-模同态
\begin{equation}
  B_k\overline C[1]\longrightarrow\overline C[1]
  \tag{3.2.13.3}\label{eq:3-2-13-3}
\end{equation}
诱导。换言之，我们假定 $\overline{\mathfrak m}_k$ 不含 $e,e^{-1}$，并仍用
$\overline{\mathfrak m}_k$ 表示~\eqref{eq:3-2-13-3} 中的 $R$-模同态。

现在直和 $B(C[1]):=\bigoplus_kB_k(C[1])$ 具有分次余代数结构。我们把
$B(C[1])$ 看作余代数，并在其上构造余导子。余积 $\Delta$ 定义为
\begin{equation}
  \Delta(x_1\otimes\cdots\otimes x_n)
  =\sum_{k=0}^{n}
  (x_1\otimes\cdots\otimes x_k)
  \otimes(x_{k+1}\otimes\cdots\otimes x_n).
  \tag{3.2.14}\label{eq:3-2-14}
\end{equation}
可以把 $\mathfrak m_k$ 唯一延拓为余导子\liftIndex{I092}
\[
  \widehat{\mathfrak m}_k\colon
  \bigoplus_nB_n(C[1])
  \longrightarrow
  \bigoplus_nB_{n-k+1}(C[1]),
\]
其定义为
\begin{equation}
\begin{aligned}
  &\widehat{\mathfrak m}_k(x_1\otimes\cdots\otimes x_n)\\
  &\quad={}
  \sum_{\ell=1}^{n-k+1}
  (-1)^{\deg x_1+\cdots+\deg x_{\ell-1}+\ell-1}
  x_1\otimes\cdots\otimes
  \mathfrak m_k(x_\ell,\ldots,x_{\ell+k-1})
  \otimes\cdots\otimes x_n .
\end{aligned}
\tag{3.2.15}\label{eq:3-2-15}
\end{equation}
这里 $k\leq n$；当 $k>n$ 时，令 $\widehat{\mathfrak m}_k=0$。在
$k=0$ 的情形，把右端的 $\mathfrak m_k$ 换成 $\mathfrak m_0(1)$；也就是说，
\[
\begin{aligned}
  &\widehat{\mathfrak m}_0(x_1\otimes\cdots\otimes x_n)\\
  &\quad={}
  \sum_{\ell=1}^{n+1}
  (-1)^{\deg x_1+\cdots+\deg x_{\ell-1}+\ell-1}
  x_1\otimes\cdots\otimes x_{\ell-1}\otimes
  \mathfrak m_0(1)\otimes x_\ell\otimes\cdots\otimes x_n.
\end{aligned}
\]

下面考虑 $B(C[1])$ 的完备化 $\widehat B(C[1])$。在 $B_k(C[1])$ 上定义滤过
\[
  F^\lambda B_k(C[1])
  =\bigcup_{\lambda_1+\cdots+\lambda_k\geq\lambda}
  \bigl(F^{\lambda_1}C^{m_1}\otimes\cdots\otimes
  F^{\lambda_k}C^{m_k}\bigr).
\]
这里的分次指标 $m_1,\ldots,m_k$ 按各分次分量逐一理解。
令 $\widehat B_k(C[1])$ 为 $B_k(C[1])$ 关于这个滤过的完备化；仍称该滤过为
能量滤过\liftIndex{I081}\liftIndex{I093}。$B(C[1])$ 上还有由
\[
  \bigoplus_{k\geq k_0}B_k(C[1]),
  \qquad k_0=1,2,\ldots,
\]
给出的另一滤过，称为\emph{数目滤过}\liftIndex{I094}。

\begin{definition}{3.2.16}\label{def:3-2-16}
$\widehat B(C[1])$ 是所有形式和 $\sum_kx_k$ 构成的集合，其中
$x_k\in\widehat B_k(C[1])$，并且存在 $\lambda_k$ 使
\[
  x_k\in F^{\lambda_k}\widehat B_k(C[1]),
  \qquad \lambda_k\longrightarrow\infty\quad(k\longrightarrow\infty).
\]
\end{definition}

定义
\[
  \widehat d\colon\widehat B(C[1])\longrightarrow\widehat B(C[1]),
  \qquad
  \widehat d=\sum_{k=0}^{\infty}\widehat{\mathfrak m}_k.
\]

\begin{lemma}{3.2.17}\label{lem:3-2-17}
若 (3.2.12) 成立，则 $\widehat d$ 是从 $\widehat B(C[1])$ 到自身的
良定义映射。
\end{lemma}

\begin{proof}
令 $\sum_kx_k$ 如定义 3.2.16。$\widehat d\sum_kx_k$ 的
$\widehat B_k(C[1])$-分量是
$\sum_\ell\widehat{\mathfrak m}_\ell(x_{k+\ell-1})$。由 (3.2.12) 和定义
3.2.16，
\[
  \widehat{\mathfrak m}_\ell(x_{k+\ell-1})
  \in F^{\lambda_{k+\ell-1}}\widehat B_k(C[1]).
\]
因此，$\widehat B_k(C[1])$ 的完备性蕴含
$\widehat d\sum_kx_k$ 收敛，引理得证。
\end{proof}

现在对元素 $e\in C^0=C[1]^{-1}$ 引入下述条件；另可比较
\textbf{[Fuk97II]} 第 11 节。

\begin{statement}{条件}{3.2.18}\label{cond:3-2-18}
\begin{align}
  \mathfrak m_{k+1}(x_1,\ldots,e,\ldots,x_k)&=0,
  &&k\geq2\text{ 或 }k=0,
  \tag{3.2.19.1}\label{eq:3-2-19-1}\\
  \mathfrak m_2(e,x)&=(-1)^{\deg x}\mathfrak m_2(x,e)=x.
  \tag{3.2.19.2}\label{eq:3-2-19-2}
\end{align}
\end{statement}

\pagebreak
\begin{definition}{3.2.20}\label{def:3-2-20}
\begin{enumerate}[label=\textup{(\arabic*)},leftmargin=2.8em]
  \item $C$ 上的一个\emph{滤过 $A_\infty$ 代数结构}
  \liftIndex{I002}\liftIndex{I082}，是满足
  (3.2.12.1)--(3.2.12.7)、(3.2.13.1)--(3.2.13.3) 以及
  $\widehat d\circ\widehat d=0$ 的一族
  $\mathfrak m=\{\mathfrak m_k\}_{k\geq0}$。称
  $\widehat B(C[1])$ 为与 $A_\infty$ 代数 $(C,\mathfrak m)$ 相伴的
  （完备）bar 复形\liftIndex{I095}。

  若 $(C,\mathfrak m)$ 是滤过 $A_\infty$ 代数，则
  (3.2.13) 中的
  $(\overline C,\overline{\mathfrak m}_k)$ 是 $A_\infty$ 代数。称它为
  $(C,\mathfrak m_k)$ 的 $R$-约化\liftIndex{I084}。本书常把
  $(C,\mathfrak m_k)$ 等的 $R$-约化直接记作
  $(\overline C,\overline{\mathfrak m}_k)$，不再另行说明。

  \item 若滤过 $A_\infty$ 代数 $(C,\mathfrak m)$ 满足
  $\mathfrak m_0=0$，就称它是\emph{严格的}\liftIndex{I085}。

  \item 若一个（滤过）$A_\infty$ 代数含有次数为 $0$（移位前）的元素
  $e$，并且 $e$ 满足条件 3.2.18，就称该代数为\emph{带幺
  （滤过）$A_\infty$ 代数}，称 $e$ 为它的\emph{幺元}
  \liftIndex{I086}；有时也简称\emph{含幺滤过 $A_\infty$ 代数}。
  对含幺滤过 $A_\infty$ 代数 $(C,\mathfrak m)$，置
  \[
    C_{\mathrm{red}}=C/\Lambda_{0,\mathrm{nov}}[e],
  \]
  并且本书始终假定存在分裂
  \[
    C=C_{\mathrm{red}}\oplus\Lambda_{0,\mathrm{nov}}[e].
  \]
  对 $R$ 上的无滤过 $A_\infty$ 代数
  $(\overline C,\overline{\mathfrak m})$，也以同样方式定义幺元并假定上述
  分裂存在。关于 $C_{\mathrm{red}}$ 的一些讨论，见第 4.4.4、4.4.7 小节。
\end{enumerate}
\end{definition}

\begin{remark}{3.2.21}\label{rem:3-2-21}
\begin{enumerate}[label=\textup{(\arabic*)},leftmargin=2.8em]
  \item 方程 $\widehat d\circ\widehat d=0$ 在各个 $\mathfrak m_k$ 之间产生
  无穷多个关系。例如，最初几个关系为
  \[
  \begin{aligned}
    &\mathfrak m_1(\mathfrak m_0(1))=0,\\
    &\mathfrak m_2(\mathfrak m_0(1),x)
      +(-1)^{\deg x+1}\mathfrak m_2(x,\mathfrak m_0(1))\\
    &\qquad{}+\mathfrak m_1(\mathfrak m_1(x))=0,\\
    &\mathfrak m_3(\mathfrak m_0(1),x,y)
      +(-1)^{\deg x+1}\mathfrak m_3(x,\mathfrak m_0(1),y)\\
    &\qquad{}+(-1)^{\deg x+\deg y+2}
      \mathfrak m_3(x,y,\mathfrak m_0(1))\\
    &\qquad{}+\mathfrak m_2(\mathfrak m_1(x),y)
      +(-1)^{\deg x+1}\mathfrak m_2(x,\mathfrak m_1(y))\\
    &\qquad{}+\mathfrak m_1(\mathfrak m_2(x,y))=0,
  \end{aligned}
  \]
  等等。一般地，容易证明 $\widehat d\circ\widehat d=0$ 等价于：对每个
  $k$，都有
  \begin{equation}
  \begin{aligned}
    &\sum_{k_1+k_2=k+1}\sum_i
    (-1)^{\deg x_1+\cdots+\deg x_{i-1}+i-1}\\
    &\qquad\mathfrak m_{k_1}
    \bigl(x_1,\ldots,
      \mathfrak m_{k_2}(x_i,\ldots,x_{i+k_2-1}),
      \ldots,x_k\bigr)=0.
  \end{aligned}
  \tag{3.2.22}\label{eq:3-2-22}
  \end{equation}
  称这个恒等式为 $A_\infty$ 公式或 $A_\infty$ 关系。若
  $\mathfrak m_0=0$，则 $\mathfrak m_1\mathfrak m_1=0$；此时
  $\mathfrak m_1$ 起上边界算子的作用。

  \item 假设 $\mathfrak m_0=0$ 且
  $\mathfrak m_3=\mathfrak m_4=\cdots=0$。那么
  $\widehat d\circ\widehat d=0$ 蕴含
  \begin{equation}
    \mathfrak m_2(\mathfrak m_2(x,y),z)
    +(-1)^{\deg x+1}\mathfrak m_2(x,\mathfrak m_2(y,z))=0.
    \tag{3.2.23}\label{eq:3-2-23}
  \end{equation}
  式~\eqref{eq:3-2-23} 与~\eqref{eq:3-2-5} 以及 D.G.A. 中乘积的结合律相容。

  \item 置\liftIndexRangeStart{I097}
  \begin{equation}
    \mathfrak l_k(x_1,\ldots,x_k)
    =\frac1{k!}\sum_{\sigma\in S_k}(-1)^{\epsilon(\sigma)}
    \mathfrak m_k(x_{\sigma(1)},\ldots,x_{\sigma(k)}),
    \tag{3.2.24}\label{eq:3-2-24}
  \end{equation}
  其中
  \begin{equation}
    \epsilon(\sigma)
    =\sum_{i<j;\,\sigma(i)>\sigma(j)}
    (\deg x_i+1)(\deg x_j+1).
    \tag{3.2.25}\label{eq:3-2-25}
  \end{equation}
  这可以看作对算子 $\mathfrak m_k$ 的对称化
  \liftIndex{I096}。$\mathfrak m_k$ 的 $A_\infty$ 公式在算子
  $\mathfrak l_k$ 之间导出若干关系，而这些关系定义一个
  $L_\infty$ 结构\liftIndexRangeEnd{I097}。关于 $A_\infty$ 结构的对称化，
  见第 A3 节以及例如 \textbf{[LaMa95]}。有些文献采用的符号约定与本书略有不同。

  \item 第 3.5 节将构造一个与拉格朗日子流形关联的滤过 $A_\infty$ 代数。
  它没有定义 3.2.20 (3) 意义下的幺元；但第 7.3 节将证明它有同伦幺元。
  我们将在第 3.3 节定义 $A_\infty$ 代数中同伦幺元的概念。

  \item 注意，完备化 $\widehat B(C[1])$ 并不是通常意义下的余代数，但存在运算
  \[
    \Delta\colon\widehat B(C[1])
    \longrightarrow
    \widehat B(C[1])\widehat\otimes\widehat B(C[1]),
  \]
  满足
  \[
    (\Delta\widehat\otimes1)\circ\Delta
    =(1\widehat\otimes\Delta)\circ\Delta.
  \]
  这里以及下文中，$\widehat\otimes$ 表示
  $\Lambda_{0,\mathrm{nov}}$ 上张量积的完备化。称这样的代数结构为
  \emph{形式余代数}。这一概念也出现在形式群理论中
  \textbf{[Die73]}。我们的 $\widehat d$ 是这个形式余代数的余导子。
\end{enumerate}
\end{remark}

下面定义有隙条件。令
$G\subset\mathbb R_{\geq0}\times2\mathbb Z$ 为子幺半群，并把它写成
\[
  G=\{\beta=(\lambda(\beta),\mu(\beta))
  \in\mathbb R_{\geq0}\times2\mathbb Z\mid\beta\in G\}.
\]

\begin{definition}{3.2.26}\label{def:3-2-26}
假设 $G$ 满足条件 3.1.6。若 $\mathfrak m_k$ 有分解
\[
  \mathfrak m_k
  =\sum_{\beta\in G}T^{\lambda(\beta)}e^{\mu(\beta)/2}
  \mathfrak m_{k,\beta},
\]
其中对 $k=0,1,2,\ldots$ 及
$\beta=(\lambda(\beta),\mu(\beta))\in G$，
\[
  \mathfrak m_{k,\beta}\colon
  B_k\overline C[1]\longrightarrow\overline C[1]
\]
是一族 $R$-模同态，就称滤过 $A_\infty$ 代数 $(C,\mathfrak m)$ 是
\emph{$G$-有隙的}\liftIndex{I067}\liftIndex{I083}。这里把
$C$ 与 $\overline C\otimes_R\Lambda_{0,\mathrm{nov}}$ 等同。若对某个 $G$，
$(C,\mathfrak m)$ 是 $G$-有隙的，就简称它\emph{有隙}。
\end{definition}

注意，(3.1.7.2) 蕴含 (3.2.13.3)。在构造谱序列时，以及对滤过作各种归纳构造
时，都要使用有隙条件；见第 4、5、6 章。

下面定义滤过 $A_\infty$ 同态。令
$(C_i,\mathfrak m^{\,i})$（$i=1,2$）为环
$\Lambda_{0,\mathrm{nov}}$ 上的滤过 $A_\infty$ 代数。对
$k=0,1,2,\ldots$，考虑次数为 $0$ 的 $\Lambda_{0,\mathrm{nov}}$-模同态族
\[
  f_k\colon B_k(C_1[1])\longrightarrow C_2[1],
\]
满足
\begin{align}
  f_k\bigl(F^\lambda B_k(C_1[1])\bigr)
  &\subset F^\lambda C_2[1],
  \tag{3.2.27.1}\label{eq:3-2-27-1}\\
  f_0(1)&\in F^{\lambda'}C_2[1]
  \qquad\text{对某个 }\lambda'>0.
  \tag{3.2.27.2}\label{eq:3-2-27-2}
\end{align}
注意 $f_0\colon\Lambda_{0,\mathrm{nov}}\to C_2[1]$，这里并不假定
$f_0(1)=0$。这些映射诱导
\[
  \widehat f\colon\widehat B(C_1[1])\longrightarrow\widehat B(C_2[1]),
\]
它在 $B_k(C_1[1])$ 上由下式给出：
\begin{equation}
\begin{aligned}
  &\widehat f(x_1\otimes\cdots\otimes x_k)\\
  &\quad={}
  \sum_{0\leq k_1\leq\cdots\leq k_n\leq k}
  f_{k_1}(x_1,\ldots,x_{k_1})\otimes\cdots\\
  &\qquad\quad{}\otimes
  f_{k_{i+1}-k_i}(x_{k_i+1},\ldots,x_{k_{i+1}})\otimes\cdots\\
  &\qquad\quad{}\otimes
  f_{k-k_n}(x_{k_n+1},\ldots,x_k).
\end{aligned}
\tag{3.2.28}\label{eq:3-2-28}
\end{equation}
条件 (3.2.27) 蕴含 $\widehat f$ 收敛。当 $f_0$ 出现在
\eqref{eq:3-2-28} 右端时，以 $f_0(1)$ 代入。特别地，按本书的指数函数记号，
\[
  \widehat f(1)
  =1+f_0(1)+f_0(1)\otimes f_0(1)+\cdots=e^{f_0(1)}.
\]
另见下面的 (3.6.1)。由 (3.2.27.2)，上式右端收敛。容易看出
\[
  \widehat f\colon\widehat B(C_1[1])\longrightarrow\widehat B(C_2[1])
\]
是形式余代数同态；另见引理 3.6.2。

\begin{definition}{3.2.29}\label{def:3-2-29}
令 $f_k\colon B_k(C_1[1])\to C_2[1]$ 如上满足 (3.2.27)。我们假定
$R[e,e^{-1}]$-模同态\translatornote{原书在这里的两个映射中都把
源、目标的下标 $1,2$ 漏去；但本定义开头已指定
$f_k\colon B_k(C_1[1])\to C_2[1]$，故在约化后补为
$B_k\overline C_1[1]\to\overline C_2[1]$。}
\[
  f_k\bmod\Lambda^+_{0,\mathrm{nov}}\colon
  B_k\overline C_1[1]\otimes_RR[e,e^{-1}]
  \longrightarrow
  \overline C_2[1]\otimes_RR[e,e^{-1}]
\]
由某个 $R$-模同态
\[
  \overline f_k\colon
  B_k\overline C_1[1]\longrightarrow\overline C_2[1]
\]
诱导。
\begin{enumerate}[label=\textup{(\arabic*)},leftmargin=2.8em]
  \item 若
  $\widehat f\circ\widehat d^{\,1}
  =\widehat d^{\,2}\circ\widehat f$，就称
  $f=\{f_k\}_{k\geq0}$ 为从 $C_1$ 到 $C_2$ 的
  \emph{滤过 $A_\infty$ 同态}\liftIndex{I089}，简记为
  $f\colon C_1\to C_2$。

  \item 若 $f_0=0$，就称滤过 $A_\infty$ 同态 $f$ 是
  \emph{严格的}\liftIndex{I091}。

  \item 令 $G\subset\mathbb R_{\geq0}\times2\mathbb Z$ 为满足条件 3.1.6
  的离散子幺半群，并令 $C_i$（$i=1,2$）为有隙滤过 $A_\infty$ 代数。
  若存在映射
  $f_{k,\beta}\colon B_k\overline C_1[1]\to\overline C_2[1]$
  （$k=0,1,2,\ldots$，$\beta\in G$），使得
  \begin{equation}
    f_k=\sum_{\beta\in G}f_{k,\beta}
    T^{\lambda(\beta)}e^{\mu(\beta)/2},
    \tag{3.2.30}\label{eq:3-2-30}
  \end{equation}
  就称 $f$ 是\emph{$G$-有隙的}\liftIndex{I090}。
\end{enumerate}
\end{definition}

如果 $G'$ 是另一个满足条件 3.1.6 并且包含 $G$ 的离散子幺半群，那么每个
$G$-有隙滤过 $A_\infty$ 同态也都是 $G'$-有隙的。不明确指定子幺半群
$G$ 时，只说 $f$ 是有隙的。

设 $f_k^{\,i}\colon B_k(C_i[1])\to C_{i+1}[1]$（$i=1,2$）定义了滤过
$A_\infty$ 同态。

\begin{definition}{3.2.31}\label{def:3-2-31}
$f^{\,1}$ 与 $f^{\,2}$ 的复合
$f^{\,2}\circ f^{\,1}=\{(f^{\,2}\circ f^{\,1})_k\}$
\liftIndex{I088}由下式定义：
\[
\begin{aligned}
  &(f^{\,2}\circ f^{\,1})_k(x_1,\ldots,x_k)\\
  &\quad={}
  \sum_m\sum_{k_1+\cdots+k_m=k}
  f_m^{\,2}\bigl(
    f_{k_1}^{\,1}(x_1,\ldots,x_{k_1}),\ldots,
    f_{k_m}^{\,1}(x_{k-k_m+1},\ldots,x_k)
  \bigr).
\end{aligned}
\]
\end{definition}

\begin{lemma}{3.2.32}\label{lem:3-2-32}
$f^{\,2}\circ f^{\,1}$ 定义了从 $C_1$ 到 $C_3$ 的滤过
$A_\infty$ 同态。
\end{lemma}

\begin{proof}
注意到 $\widehat f(1)=e^{f_0(1)}$，证明就与引理 3.2.9 的证明相同。
\end{proof}

令 $(C_i,\mathfrak m^{\,i})$（$i=1,2$）为
$\Lambda_{0,\mathrm{nov}}$ 上的滤过 $A_\infty$ 代数，并令
$f\colon(C_1,\mathfrak m^{\,1})\to(C_2,\mathfrak m^{\,2})$ 为滤过
$A_\infty$ 同态。那么 $f$ 自然诱导 $A_\infty$ 同态
\[
  \overline f\colon
  (\overline C_1,\overline{\mathfrak m}^{\,1})
  \longrightarrow
  (\overline C_2,\overline{\mathfrak m}^{\,2}),
\]
其中 $(\overline C_i,\overline{\mathfrak m}^{\,i})$ 是
$(C_i,\mathfrak m^{\,i})$ 的 $R$-约化。

\begin{definition}{3.2.33}\label{def:3-2-33}
令 $(C_i,\mathfrak m^{\,i})$（$i=1,2$）为
$\Lambda_{0,\mathrm{nov}}$ 上的滤过 $A_\infty$ 代数，并且
$\overline{\mathfrak m}^{\,i}_0=0$。若滤过 $A_\infty$ 同态
\[
  f\colon(C_1,\mathfrak m^{\,1})
  \longrightarrow(C_2,\mathfrak m^{\,2})
\]
所诱导的 $A_\infty$ 同态
$\overline f\colon(\overline C_1,\overline{\mathfrak m}^{\,1})
\to(\overline C_2,\overline{\mathfrak m}^{\,2})$ 是定义 3.2.10 意义下的
弱同伦等价，就称 $f$ 是\emph{弱同伦等价}\liftIndex{I098}。进一步，若
$(C_i,\mathfrak m^{\,i})$ 与 $f$ 都是 $G$-有隙的，就称它为
\emph{$G$-有隙弱同伦等价}。
\end{definition}

\begin{definition}{3.2.34}\label{def:3-2-34}
对滤过 $A_\infty$ 代数 $(C,\mathfrak m)$，若其 $R$-约化与某个滤过
$A_\infty$ 代数 $(C',\mathfrak m')$ 的 $R$-约化之间存在弱同伦等价
\[
  \overline f\colon(\overline C,\overline{\mathfrak m})
  \longrightarrow(\overline C',\overline{\mathfrak m}'),
\]
就称 $(C',\mathfrak m')$ 是 $(C,\mathfrak m)$ 的一个
\emph{$A_\infty$ 形变}\liftIndex{I080}。类似地，若无滤过
$A_\infty$ 代数 $(\overline C,\overline{\mathfrak m})$ 与滤过
$A_\infty$ 代数 $(C',\mathfrak m')$ 的 $R$-约化之间存在弱同伦等价
\[
  \overline f\colon(\overline C,\overline{\mathfrak m})
  \longrightarrow(\overline C',\overline{\mathfrak m}'),
\]
也称 $(C',\mathfrak m')$ 是
$(\overline C,\overline{\mathfrak m})$ 的一个 $A_\infty$ 形变。
\end{definition}
